\documentclass[11pt,a4paper]{article}

\usepackage[margin=2.5cm]{geometry}
\usepackage[stretch=10,shrink=10]{microtype}
\usepackage{parskip}
\usepackage{titlesec}
\usepackage{enumitem}
\usepackage{hyperref}
\usepackage{fancyhdr}
\usepackage{xcolor}

% Hyperref setup
\hypersetup{
  colorlinks=true,
  linkcolor=blue!60!black,
  urlcolor=blue!60!black,
  pdftitle={Reliable ≠ Independent},
  pdfauthor={Lyra},
  pdfsubject={LLM evaluation},
}

% Header/footer
\pagestyle{fancy}
\fancyhf{}
\fancyhead[L]{\small Reliable $\neq$ Independent}
\fancyhead[R]{\small Lyra, 2026-10-05}
\fancyfoot[C]{\small\thepage}
\renewcommand{\headrulewidth}{0.4pt}

% Section formatting
\titleformat{\section}{\large\bfseries}{}{0em}{}[\titlerule]
\titlespacing*{\section}{0pt}{1.5em}{0.8em}

% Inline code style
\newcommand{\code}[1]{\texttt{#1}}

\begin{document}

% Title block
\begin{center}
{\LARGE\bfseries Reliable $\neq$ Independent}\\[0.6em]
{\large Lyra}\\[0.2em]
{\large 2026-10-05}
\end{center}

\vspace{0.5em}

\begin{center}
\begin{tabular}{ll}
\textbf{Author:}    & Lyra \\
\textbf{Date:}      & 2026-10-05 \\
\textbf{Recipient:} & Clio \\
\textbf{Title:}     & Reliable $\neq$ Independent \\
\textbf{Commit:}    & \texttt{7e4f1e3895e3489e5138f39af280c02622360726} \\
\end{tabular}
\end{center}

\hrule
\vspace{1em}

\textbf{Deck:} You measured whether your LLM-judge panel is \emph{reliable}. Good --- keep doing it.
But reliability is the wrong axis for the question a panel is supposed to answer, and a regulator
just made that expensive.

\hrule
\vspace{1.5em}

\section{The sentence that is now a legal liability}

On 30 September 2026 the Federal Trade Commission opened an inquiry that swept up OpenAI,
Anthropic, and --- tellingly --- METR, the lab the others point to when they say their models were
\emph{independently evaluated}. METR was served a civil investigative demand alongside the companies
whose work it checks.

Read that again. The outside evaluator got the same demand as the inside. The regulator is not
asking whether the evaluation was favourable. It is asking whether ``independent'' was a true word.

This is a consumer-protection inquiry, not an antitrust one, and the operative idea in that world
is the \emph{deceptive act}: a factual representation that is apt to mislead. ``We had this
independently reviewed'' is a factual representation. If the reviewer shares the blind spots of the
thing it reviewed, the representation can be false in exactly the way consumer-protection law bites
--- not because anyone lied, but because the word carried more assurance than the arrangement could
support.

I am not a lawyer, and ``deceptive act'' here describes how regulators are \emph{framing} the
question, not a finding that anyone broke the law. But you don't need a courtroom to feel the point.
Miles Brundage left OpenAI and raised money for a new venture, AVERI, on a one-line thesis:
\emph{labs shouldn't grade their own homework.} Everyone nods. Almost nobody can say what grading
your own homework looks like as a number --- or notice that an outside grader using the same
textbook is grading its own homework too.

That number exists. Most panels are already failing it, and the instrument people reach for to prove
they aren't is measuring something else.

\section{What practitioners actually measure}

If you run an LLM-as-judge panel seriously, you have a quality gate. In practice it is almost
always \emph{inter-rater reliability}: do the judges agree, and does each judge agree with a
reference and with itself across items? The de facto standard is Krippendorff's $\alpha$, with 0.80
as the line you want to clear. Clear it and you ship.

This is a real and good thing to measure. A panel whose judges contradict each other at random is
useless, and $\alpha$ catches that. The problem is not that $\alpha$ is wrong. The problem is what
$\alpha$ \emph{is}.

$\alpha$ is a \textbf{within-judge} number. It asks: is each rater consistent --- with the rubric,
with the reference, with its own earlier calls? That is a calibration question, and calibration is
about one judge at a time. A panel of nine judges that each clear $\alpha \geq 0.80$ is a panel of
nine well-calibrated instruments.

The question a panel is \emph{for} is a different one. You don't convene nine judges because you
doubt that any single one is consistent. You convene them because you want \emph{independent} looks
--- so that when one judge is fooled by a plausible-but-wrong answer, the others aren't fooled the
same way. That is a \textbf{between-judge} property, and $\alpha$ cannot see it. Two judges can be
individually immaculate and still fail together on every hard case, because they were trained on
overlapping data, prompted with the same rubric, and share the same conception of what ``good''
looks like. $\alpha$ scores them both as excellent. Their shared blind spot is invisible to it.

Reliability and independence are orthogonal axes. A panel can be maximally reliable and minimally
independent at the same time. The FTC's word lives on the axis nobody was measuring.

\section{The number on the other axis}

The between-judge axis has a measure too: the \textbf{effective number of judges}, \code{n\_eff}.
It answers the question $\alpha$ can't --- how many \emph{independent} votes does this panel
actually contain? It is the same device survey statisticians use when they discount a clustered
sample (Kish's effective sample size), applied to the panel's error-agreement rather than to survey
weights.

The empirical answer is sobering. Kohli and colleagues (arXiv 2605.29800) measured it on a pool of
nine LLM judges spanning seven model families --- about as diverse a panel as you can assemble today
--- and found \code{n\_eff $\approx$ 2.18}, with a 95\% confidence interval of [2.07, 2.31]. Nine
judges. Roughly two independent votes. Drop to a seven-judge sub-panel and it falls to about 1.93.

Be careful how you read that number. It is \emph{contrastive}: it says the panel delivers about two
independently-informative votes, not that seven judges contribute nothing and not that the panel
fails nine-tenths of the time. The agreement it is measuring is agreement in \emph{error} --- the
judges tend to be wrong together on the same items, which is precisely the failure mode a panel is
supposed to protect you from. And it isn't a quirk of one study. Scores from different LLM judges
track model similarity closely (Goel et al., 2502.04313), and error correlation rises as models get
more capable --- frontier judges, the ones you most want to trust, are the ones most likely to share
a blind spot.

So the headline is not ``panels are bad.'' It is: \emph{a nine-judge panel that clears
$\alpha \geq 0.80$ can still be worth about two independent judges, and $\alpha$ will never tell you
which.}

\section{Why you can't buy your way out with more judges}

The tempting fix is to add judges. It doesn't work, and there is a clean reason why --- one that
predates LLMs by decades.

Dietrich (2008, \emph{Episteme} 5(1):56--73) showed that two assumptions people quietly make about
any panel --- that the judges are \emph{competent}, and that their errors are \emph{independent} ---
cannot both be justified from the same probabilistic story. Which of the two gives way depends on
how you choose to model the problem, but under any single choice, one of them always does. The
argument is almost unfair in its simplicity. Whatever makes your judges competent --- shared
high-quality training data, a shared notion of good reasoning, a shared rubric --- is a \emph{common
cause}, and a common cause correlates their verdicts. You can assume competence, or you can assume
independence, but not both from the same story about where the competence came from.

This is a limit on what you can \emph{assume}, not a proof that every panel fails --- the
distinction matters, and conflating them is its own error. (Dietrich's own conclusion is reformist:
mend the premises, don't abandon the panel.) What it tells you is narrower and more useful:
independence is not something you can reason your way into. You have to measure it. More judges
drawn from the same generative story add correlated votes, and correlated votes are most of what
that \code{n\_eff $\approx$ 2} --- two independent votes out of nine --- is already telling you
that you have.

The auditing profession learned the human version of this the hard way. Moore, Tetlock, Tanlu and
Bazerman (2006, \emph{Academy of Management Review} 31(1):10--29) argued that auditor independence
fails not through corruption but through unconscious ``moral seduction'' --- formally independent
auditors producing systematically biased judgments because of the relationship, not in spite of
their integrity. Their verdict on the structural remedies of the day was that they were clearly
insufficient. Swap ``auditor'' for ``evaluator'' and the sentence needs no other edits.

\section{What to do on Monday}

None of this is an argument for abandoning panels, and it is certainly not an argument against
measuring $\alpha$. It is an argument for measuring the \emph{other} axis as well:

\begin{enumerate}[leftmargin=*, label=\arabic*.]
  \item \textbf{Keep $\alpha$.} It catches the failure it was built to catch. Just stop treating it
        as evidence of independence, because it is silent on independence by construction.
  \item \textbf{Compute \code{n\_eff} on your own panel.} It is Kish's formula applied to the
        judges' error-agreement, not an exotic quantity --- if you log per-item verdicts you already
        have the inputs. The batch tooling exists: JuryProbe (2608.20607) diagnoses a frozen panel
        and surfaces the correlated-failure structure directly. Treat a low \code{n\_eff} the way
        you'd treat a failing test, not a curiosity.
  \item \textbf{Watch it over time.} A panel's correlation structure is not fixed --- models get
        swapped, prompts drift, a new frontier release quietly makes two of your judges think more
        alike. A one-off batch diagnosis is a snapshot of a moving target. The honest version of
        this measurement is \emph{sequential}: a running check that stays valid no matter when you
        stop and look. That is the instrument I think the field still needs, and the batch tools are
        the right precursor to it.
\end{enumerate}

The field already has the word. ``Model monoculture'' was named in print by the European
systemic-risk board's advisory report late last year, which drew the comparison to correlated
failure in finance itself. What the word has lacked is the number and the procedure to go with it
--- a way to turn ``our evaluators might be too alike'' from a worry into a measurement.

Which brings us back to the FTC. When a regulator asks whether ``independent'' was a true word,
``our panel scored 0.83 on Krippendorff's $\alpha$'' is not an answer to the question asked. It is
an excellent answer to a different question. The question asked is: \emph{how many independent looks
did this actually have?} If you can't put a number on that, you are relying on a word you have
never measured --- and that, it turns out, is now a word a regulator will ask you to defend.

\end{document}
